\chapter{Convincing a Paranoid Kernel That a 77-Digit Number Is Prime}
\label{ch:prime}

\section{The problem nobody warns you about}

Chapter~\ref{ch:modular} ended with a quiet dependency: everything ---
division, field structure, the whole elliptic curve --- rests on
$p = 2^{255}-19$ \emph{being prime}. In Lean, that is a proposition like any
other, and it must be \emph{proved}:

\begin{lstlisting}[language=Lean]
theorem p_prime : Nat.Prime (2^255 - 19) := ?
\end{lstlisting}

Your Chapter~\ref{ch:automation} instincts say \lean{decide}: primality is
decidable --- just try dividing. But trial division tests divisors up to
$\sqrt{p} \approx 2^{127}$. At a billion billion divisions per second, that
is about $10^{12}$ ages of the universe. The kernel, which happily
\emph{re-executes} every computation you feed it, cannot afford this one. And
mathematicians clearly believe this number is prime --- so how does
\emph{anyone} know?

\section{Certificates: the deep idea hiding here}

The answer reorganizes how you think about computation. \emph{Finding} a
fact and \emph{checking} a fact can have wildly different costs. What we need
is a \textbf{certificate}: a piece of data, possibly expensive to discover,
that makes the fact \emph{cheap to verify}.

You have met certificates before without the name. A composite number's
certificate is a factor: finding a factor of a 77-digit number may be hard,
but checking $n = a \cdot b$ is one multiplication. The beautiful surprise
--- Pratt's theorem, 1975 --- is that \emph{primality} has certificates too:

\begin{bigidea}
\textbf{Pratt certificate.} To certify that $p$ is prime, exhibit a
\emph{witness} $w$ such that
\[
w^{p-1} \equiv 1 \pmod p
\qquad\text{and}\qquad
w^{(p-1)/q} \not\equiv 1 \pmod p
\ \text{ for every prime factor } q \text{ of } p-1 .
\]
Such a $w$ generates all $p-1$ nonzero residues, which forces $\Zmod{p}$ to
have $p-1$ invertible elements --- something only a prime modulus allows.
Checking the certificate costs a handful of modular exponentiations
(milliseconds, by fast squaring), \emph{plus recursively certifying the
prime factors $q$} --- each much smaller, so the recursion collapses fast.
\end{bigidea}

Concretely for our hero: $p - 1 = 2^{255} - 20$ factors as
\[
p - 1 \;=\; 2^{2} \cdot 3 \cdot 65147 \cdot Q,
\]
where $Q$ is a 71-digit prime with its own (short) certificate, and
$65147$ recurses one more level: $65146 = 2 \cdot 32573$ with $32573$
prime. The full certificate for $p$ is a small tree of witnesses and
factorizations --- a few hundred bytes of data standing behind a 77-digit
claim:

\begin{center}
\begin{tikzpicture}[
  lvl/.style={draw=ink2,thick,rounded corners=2pt,fill=white,align=center,font=\small},
  edge/.style={-{Stealth},ink2,thick}
]
  \node[lvl,fill=accentsoft] (p) at (0,2.7)
    {$p = 2^{255}-19$ \quad witness $w=2$\\ $p-1 = 2^2\cdot 3\cdot 65147\cdot Q$};
  \node[lvl] (two)   at (-5.0,0.9) {$2$: prime\\ \footnotesize (immediate)};
  \node[lvl] (three) at (-2.7,0.9) {$3$: prime\\ \footnotesize (immediate)};
  \node[lvl,fill=provensoft] (m) at (0.4,0.9)
    {$65147$ \quad witness $2$\\ \footnotesize $65146 = 2 \cdot 32573$};
  \node[lvl,fill=provensoft] (q) at (4.4,0.9)
    {$Q$ (71 digits)\\ own witness + factors};
  \node[lvl] (sub) at (0.4,-0.7) {$32573$: prime\\ \footnotesize (small cert)};
  \draw[edge] (p) -- (two); \draw[edge] (p) -- (three);
  \draw[edge] (p) -- (m);   \draw[edge] (p) -- (q);
  \draw[edge] (m) -- (sub);
  \node[font=\small\color{ink2},align=center] at (4.4,-0.7)
    {each node: milliseconds to check;\\ the whole tree: a proof};
\end{tikzpicture}
\end{center}

(Check one leaf of this tree yourself right now, no computer:
$4 \cdot 3 \cdot 65147 = 781764$, and $781764 \cdot Q$ must reproduce the
77-digit $p-1$ --- the \emph{product identity} is the easiest condition of
a certificate to audit, and auditing one leaf by hand is a good habit
before trusting a tree. The witness conditions for $w = 2$ at the top node
were re-verified computationally while writing this chapter; the
kernel-checked version of this construction, for the Pallas modulus, lives
in \code{pasta-pallas-verified}.)

\begin{worked}{checking a Pratt certificate by hand --- all of it}
Nothing builds trust in a certificate like verifying one completely, so
here is the full check for $n = 97$, witness $w = 5$ --- every modular
multiplication on this page, using the same square-and-multiply ladder
the real code uses (and which you will implement in the exercises).
First the data: $n - 1 = 96 = 2^5 \cdot 3$, so there are three
conditions: $5^{96} \equiv 1$, $5^{96/2} = 5^{48} \not\equiv 1$, and
$5^{96/3} = 5^{32} \not\equiv 1 \pmod{97}$.

Build the powers of $5$ by repeated squaring mod $97$, reducing as you go
--- each line is one two-digit multiplication and one division with
remainder, nothing more:
\[
\begin{array}{lclcl}
5^{2} &=& 25 \\
5^{4} &=& 25^2 = 625 &=& 6\cdot 97 + 43 \;\to\; 43\\
5^{8} &=& 43^2 = 1849 &=& 19\cdot 97 + 6 \;\to\; 6\\
5^{16} &=& 6^2 &=& 36\\
5^{32} &=& 36^2 = 1296 &=& 13\cdot 97 + 35 \;\to\; 35\\
\end{array}
\]
Now assemble the three exponents from these squares:
\[
5^{48} = 5^{32} \cdot 5^{16} = 35 \cdot 36 = 1260 = 12 \cdot 97 + 96
\;\to\; 96 \equiv -1 ,
\]
\[
5^{96} = (5^{48})^2 \equiv (-1)^2 = 1 . \qquad ✓
\]
Check the conditions: $5^{96} \equiv 1$ ✓; $5^{48} \equiv 96 \neq 1$ ✓;
$5^{32} \equiv 35 \neq 1$ ✓. All three hold --- and by Pratt's theorem
(whose reason the exercises make you articulate), $97$ is prime, with the
whole verification costing \emph{seven} small multiplications. Trial
division would have cost eight test divisions here --- no savings at two
digits. But the ladder's cost grows with the \emph{number of bits} (one
squaring per bit), while trial division grows with the \emph{square root
of the value} (one division per candidate) --- linear versus exponential
in the bit-length. At 77 digits that gap is the whole story, as the next
box counts.
\end{worked}

\begin{worked}{costing the real certificate for $p = 2^{255}-19$}
How much work is the full certificate check for the real prime, versus
trial division? Count it, honestly, using the tree above.

\emph{Certificate side.} One modular exponentiation with a 255-bit
exponent costs at most $254$ squarings plus at most $254$ multiplies ---
call it $\le 508$ modular multiplications, and abbreviate ``modmul.''
The top node needs: $2^{p-1}$ (one exponentiation), and one
$2^{(p-1)/q}$ for each of the four prime factors $q \in \{2, 3, 65147,
Q\}$ --- five exponentiations, $\le 2540$ modmuls. The recursion adds:
$Q$'s own node ($Q-1$ has its own small factor list; generously, another
five exponentiations at 71 digits, $\le 2350$ modmuls), the $65147$ node
(16-bit numbers --- three exponentiations of $\le 32$ modmuls, noise),
and $32573$'s (noise). Round the entire tree up to
\[
\text{certificate check} \;\lesssim\; 5{,}000 \text{ modmuls.}
\]
\emph{Trial division side.} $\sqrt{p} \approx 2^{127.5}$, and candidate
divisors (odd numbers, say) number about $2^{126.5} \approx 10^{38}$.
\emph{Ratio}:
\[
\frac{10^{38} \text{ divisions}}{5 \times 10^{3} \text{ modmuls}}
\;\approx\; 10^{34}.
\]
Thirty-four orders of magnitude --- not an optimization, a different
universe. And one more accounting worth doing: the certificate
\emph{data} is four factor entries and a handful of witnesses --- a few
hundred bytes. Someone (a computer algebra system, years of CPU time,
once, offline) paid dearly to \emph{find} the factorization of $p-1$;
every checker since pays five thousand multiplications. That asymmetry
--- expensive find, cheap check, tiny certificate --- is the shape of
every proof object the Lean kernel will ever hand you.
\end{worked}

\begin{aha}
This find/check asymmetry is one of the great ideas of computer science ---
it is the P versus NP distinction wearing work clothes, and it is the engine
of zero-knowledge proof systems (the very technology the Pasta curves serve).
Proof assistants run on it too: Lean's whole architecture --- clever tactics
\emph{finding}, dumb kernel \emph{checking} --- is the same asymmetry. A
proof \emph{is} a certificate.
\end{aha}

\section{The tempting shortcut, and why the house declines it}

Lean offers a faster \lean{decide}: the variant \lean{native_decide}
compiles the decision procedure to native machine code, runs it at full
speed, and asserts the result. With a good primality test behind it, it can
dispatch \lean{Nat.Prime p} in seconds. Case closed?

Look at what you would be trusting. Ordinary \lean{decide} produces a
computation the \emph{kernel} replays --- the ~few-thousand-line paranoid
core remains the only thing you trust. \lean{native_decide} instead makes
the theorem's truth depend on the Lean \emph{compiler}, the C toolchain
behind it, and the runtime --- hundreds of thousands of lines promoted into
your trusted base, in exchange for convenience on one theorem. Every proof
downstream of the field --- group law, scalars, signatures --- would inherit
that enlarged trust, visible forever in its axiom report
(Chapter~\ref{ch:honesty} shows you how to read those).

\begin{pitfall}
\lean{native_decide} is not ``cheating,'' and for exploratory work it is a
fine tool. The trap is \emph{silent trust inflation}: its use is invisible at
the theorem statement --- the cost appears only when someone audits the
axioms, which is exactly what most readers never do. House rule, adopted from
the projects this book accompanies: exploratory scaffolding may use it;
\textbf{no shipped certificate depends on it}. The final Pallas-modulus
primality proof in \code{pasta-pallas-verified} is a kernel-checked
Lucas/Pratt certificate for precisely this reason.
\end{pitfall}

\section{Certificates in practice: Mathlib's toolbox}

You will not hand-roll witness trees. Mathlib provides the machinery
(\lean{Nat.Prime} decision lemmas, \lean{lucas_lehmer}-style infrastructure,
and the \lean{norm_num} extension \lean{Nat.Prime} plugin) that constructs
and checks Pratt-style certificates behind a single tactic call --- while
keeping every step kernel-checked. The shape in real code:

\begin{lstlisting}[language=Lean]
theorem p_prime : Nat.Prime (2^255 - 19) := by
  norm_num  -- certificate-backed primality, kernel-checked, ~seconds
\end{lstlisting}

When the built-in route struggles (very large or awkward moduli), the
fallback is explicit: state the witness data as definitions, prove the two
Pratt conditions with \lean{norm_num}-driven modular exponentiation, and
assemble. That is exactly the structure of the Pallas certificate in the
companion repository --- worth reading now with fresh eyes:
\code{pasta-pallas-verified/verification/Proofs/Primality.lean}.

\begin{tryit}
Open \code{exercises/Ch07.lean}. Ladder: certify $97$, then $65537$ (a
Fermat prime beloved of RSA), then the ten-digit Mersenne prime $2^{31}-1$,
watching what each tool costs as the numbers grow. (Amusingly, the
find/check asymmetry bites the \emph{tactic} too: \lean{norm_num} must
\emph{find} the witness tree before the kernel checks it, and at $2^{61}-1$
the finding already takes minutes.) Finale: implement square-and-multiply
modular exponentiation yourself and check the top witness condition for
$p = 2^{255}-19$ with \lean{\#eval} --- your own hands on the certificate,
at 77 digits, in milliseconds.
\end{tryit}

\section*{Exercises}

\exercise{Verify by hand that $w = 2$ is a Pratt witness for $p = 13$:
compute $2^{12} \bmod 13$ and $2^{12/q} \bmod 13$ for each prime $q \mid 12$.
Write the full certificate tree for $13$, recursing into the factors of
$12$.}

\exercise{Why does the witness condition force primality? Sketch the
argument: if $w$ has order exactly $p-1$ in $\Zmod{p}$, then the
multiplicative structure has $p-1$ elements, which fails if $p = ab$ with
$1 < a,b < p$. (Full rigor optional; the shape is the point.)}

\exercise{(Paper) Certify $65147$ by hand-checkable data: given
$65146 = 2 \cdot 32573$ and witness $w = 2$, list the exact conditions a
checker must verify, then carry out the \emph{cheapest} one:
$2^{65146/32573} = 2^2 = 4 \not\equiv 1 \pmod{65147}$. For the remaining
two conditions, count precisely how many squarings and multiplications the
ladder needs (do not perform them). How many two-to-five-digit
multiplications, total, stand between a skeptic and certainty about
$65147$?}

\exercise{(Discussion) Bitcoin miners \emph{find} block hashes; nodes
\emph{check} them. GPS receivers \emph{check} satellite signals they could
never \emph{find}. Name two more systems built on the find/check asymmetry,
and one system that would collapse without it.}

\section*{Solutions and pathways}
\solutionsintro

\solhead{7.1}
\pathway The data first: $12 = 2^2 \cdot 3$, so two conditions beyond
$w^{12} \equiv 1$: exponents $12/2 = 6$ and $12/3 = 4$. Then power up $2$
mod $13$ by successive squaring, as the worked example did for $97$ ---
at this size you can even go linearly.
\answer Powers of $2$ mod $13$: $2^2 = 4$, $2^4 = 16 \equiv 3$,
$2^6 = 2^4 \cdot 2^2 = 3 \cdot 4 = 12 \equiv -1$, and
$2^{12} = (2^6)^2 \equiv (-1)^2 = 1$. Conditions: $2^{12} \equiv 1$ ✓;
$2^{6} \equiv 12 \neq 1$ ✓; $2^{4} \equiv 3 \neq 1$ ✓. Witness confirmed.
The full tree: node $13$ (witness $2$, $12 = 2^2 \cdot 3$) with children
$2$ (immediate) and $3$ (witness $2$: $2^2 = 4 \equiv 1 \pmod 3$, and
$2^{2/2} = 2 \not\equiv 1$ --- a two-line sub-certificate). Every claim in
the tree is now something you have personally multiplied.

\solhead{7.2}
\pathway The key concept is the \emph{order} of $w$: the least $e > 0$
with $w^e \equiv 1$. The two witness conditions pin the order exactly;
then count invertible elements two ways.
\answer The condition $w^{p-1} \equiv 1$ says the order of $w$ divides
$p-1$; the conditions $w^{(p-1)/q} \not\equiv 1$ for every prime
$q \mid p-1$ rule out every \emph{proper} divisor of $p-1$ (any proper
divisor of $p-1$ divides some $(p-1)/q$). So the order is exactly $p-1$:
the powers $w^1, w^2, \dots, w^{p-1}$ are $p-1$ \emph{distinct} elements,
all invertible mod $p$ (each has $w^{\text{something}}$ as inverse). But
if $p = ab$ with $1 < a, b < p$, the element $a$ is a zero divisor ---
$a \cdot b \equiv 0$ --- so $a$ is not invertible, and neither are its
multiples: strictly fewer than $p-1$ residues can be invertible.
Contradiction; $p$ has no such factorization. (Full rigor pins down
``distinct'' and the divisor-covering claim --- both one-liners with the
order concept in hand. The shape is: \emph{one loud element forces the
whole multiplicative structure to be as big as only a prime allows.})

\solhead{7.3}
\pathway List conditions mechanically from the definition, then count
ladder steps: an exponent of $b$ bits costs $\le b-1$ squarings plus (at
worst) $b-1$ multiplies; $65146 < 2^{16}$ and $32573 < 2^{15}$.
\answer Conditions: (i) $2^{65146} \equiv 1 \pmod{65147}$;
(ii) $2^{65146/2} = 2^{32573} \not\equiv 1$;
(iii) $2^{65146/32573} = 2^{2} = 4 \not\equiv 1$ ✓ (done --- four is
visibly not one); plus the recursive certificate for $32573$. Counting:
$65146$ has $16$ bits ($2^{16} = 65536$), so condition (i) costs $\le 15$
squarings $+ \le 15$ multiplies $= 30$; condition (ii), $15$ bits,
$\le 28$; condition (iii) was free. Sub-certificate for $32573$
($32572 = 2^2 \cdot 17 \cdot 479$): four conditions on $\le 15$-bit
exponents, $\le 4 \cdot 28 = 112$, plus leaves ($17$, $479$ ---
another $\sim 60$ generously). Total: \emph{under $250$ small
multiplications} --- an afternoon with paper, an eyeblink for a kernel,
and at the end $65147$ is not ``probably prime'' but \emph{prime}. (For
the audit-minded: you were given $32572$'s factorization here the same
way the checker is --- as certificate data. Verifying
$4 \cdot 17 \cdot 479 = 32572$ is one more hand multiplication:
$68 \cdot 479 = 32572$ ✓.)

\solhead{7.4}
\pathway Look for systems where producing an artifact is costly but a
short receipt convinces everyone --- then for the collapse case, imagine
checking costing as much as finding.
\answer (Model answers.) Two more: \emph{academic peer review of
computer-assisted proofs} --- the four-color theorem's checkers verify in
hours what took years to construct; and \emph{password hashing} ---
deriving a hash from a password is instant to check against, infeasible
to invert. Others students propose: sudoku (solve vs.\ check), lottery
tickets (draw vs.\ verify), digital signatures themselves (sign with
secret effort-equivalent, verify publicly). A system that would collapse
without the asymmetry: \emph{blockchain consensus} --- if verifying a
block cost as much as mining it, every node would need a mine's
electricity bill and the network could not exist. (So would mathematics
as a social enterprise: if checking a proof cost as much as finding it,
referees would be as rare as authors. In a sense, Lean is what happens
when you drive the check cost toward zero and let anyone be a referee.)

\begin{checkpoint}
You should now be able to: explain why \lean{decide} cannot prove
$2^{255}-19$ prime while a certificate can; reproduce the two Pratt witness
conditions and check them on a small prime; articulate exactly what
additional trust \lean{native_decide} would introduce and why shipped
certificates decline it; and recognize the find/check asymmetry as the
common engine of certificates, proof assistants, and the P-vs-NP question.
\end{checkpoint}
